\section{Class P --- Binary Arith --- Mod Pow}
%%% generated by the blueprint pipeline from TCSlib.Complexity.ClassP.BinaryArith.ModPow
\subsection*{Class P --- Binary Arith --- Mod Pow: Introduction}

Declarations of the Lean module \texttt{TCSlib.Complexity.ClassP.BinaryArith.ModPow}.

\subsection*{Class P --- Binary Arith --- Mod Pow: Results}

\begin{definition}[Modular exponentiation]
\lean{Complexity.BinArith.powMod}
\label{Complexity.BinArith.powMod}
\leanok
\uses{Complexity.LogProg.bitsVal, Complexity.pairFstD, Complexity.pairSndD}
For a Boolean list $z$, let $a$, $e$, and $m$ be the natural numbers encoded by,
respectively, the first component of $z$, the first component of its second component,
and the second component of its second component. Return the binary representation of
$0$ if $m=0$, and otherwise the binary representation of $a^e$ modulo $m$.
\end{definition}

\begin{theorem}[Modular Power on Encoded Inputs]
\lean{Complexity.BinArith.powMod_pairEncode}
\label{Complexity.BinArith.powMod_pairEncode}
\leanok
\uses{Complexity.BinArith.powMod, Complexity.LogProg.bitsVal_bits, Turing.pairEncode}
\difficulty{1}
For natural numbers $a,n,p$ with $p\neq 0$, encode the binary representations of $a$,
$n$, and $p$ as the nested pair
$\langle\operatorname{bits}(a),\langle\operatorname{bits}(n),\operatorname{bits}(p)\rangle\rangle$.
The binary output of the encoded modular-power operation on this input is
$\operatorname{bits}(a^n\bmod p)$.
\end{theorem}

\begin{definition}[Modular exponentiation step]
\lean{Complexity.BinArith.powStep}
\label{Complexity.BinArith.powStep}
\leanok
\uses{Complexity.BinArith.mod, Complexity.BinArith.mul, Complexity.LogProg.bitsVal, Complexity.isNilB, Complexity.pairFstD, Complexity.pairSndD, Turing.pairEncode}
For a Boolean-list state $S$, return $S$ unchanged if the first component of its second
component is empty or the natural number encoded by the first component of $S$ is zero.
Otherwise, update the encoded state using the displayed list operations: drop the first
bit of that exponent component, compute the indicated products and remainders using the
components of $S$, and choose the final remainder according to the first bit of the
exponent component.
\end{definition}

\begin{theorem}[Binary exponentiation step is polynomial-time computable]
\lean{Complexity.BinArith.polyTimeComputable_powStep}
\label{Complexity.BinArith.polyTimeComputable_powStep}
\leanok
\uses{Complexity.BinArith.mod, Complexity.BinArith.mul, Complexity.BinArith.polyTimeComputable_isZero, Complexity.BinArith.polyTimeComputable_mod, Complexity.BinArith.polyTimeComputable_mul, Complexity.BinArith.polyTimeComputable_slot2, Complexity.BinArith.polyTimeComputable_slot3, Complexity.BinArith.polyTimeComputable_slot4, Complexity.BinArith.powStep, Complexity.PolyTimeComputable, Complexity.PolyTimeComputable.comp, Complexity.PolyTimeComputable.pairEncode, Complexity.pairFstD, Complexity.pairSndD, Complexity.polyTimeComputable_headD, Complexity.polyTimeComputable_id, Complexity.polyTimeComputable_isNil, Complexity.polyTimeComputable_ite, Complexity.polyTimeComputable_or, Complexity.polyTimeComputable_pairFstD, Complexity.polyTimeComputable_tail, Turing.pairEncode}
\difficulty{4}
For a bit string $S$, define a function as follows. If the first component of the second
component of $S$ is empty, or the binary value of the first component of $S$ is zero,
return $S$. Otherwise, form a paired encoding whose first component is the first
component of $S$ and whose second component encodes, in nested pairs, the tail of the
second component of $S$, the modular product of the third component of $S$ with its
second component modulo the first component of $S$, and either the modular product of
the fourth component of $S$ with the second component of $S$ modulo the first component
of $S$ (if the first bit of the second component of $S$ is $1$), or the third component
of $S$ unchanged. Here components are obtained by successive pair decoding, failed
decodings yield the empty word, and products and remainders are computed on the
represented binary values and returned in binary. This function is polynomial-time
computable.
\end{theorem}

\begin{definition}[Modular exponentiation initialization]
\lean{Complexity.BinArith.powInit}
\label{Complexity.BinArith.powInit}
\leanok
\uses{Complexity.BinArith.mod, Complexity.pairFstD, Complexity.pairSndD, Turing.pairEncode}
For a Boolean list $z$, initialize a state by encoding the second component of the
second component of $z$, the first component of the second component of $z$, the
remainder of the first component of $z$ modulo the second component of the second
component of $z$, and the remainder of the binary representation of $1$ modulo that same
modulus.
\end{definition}

\begin{definition}[Modular exponentiation invariant]
\lean{Complexity.BinArith.PowInv}
\label{Complexity.BinArith.PowInv}
\leanok
\uses{Complexity.LogProg.bitsVal, Complexity.pairFstD, Complexity.pairSndD, Turing.pairEncode}
For Boolean lists $z$ and $S$, let $a$, $e$, and $m$ be the natural numbers encoded by
the first component of $z$, the first component of its second component, and the second
component of its second component. The state $S$ encodes lists $E$, $b$, and
$\mathrm{acc}$ in the form specified by the pair encoding, where $E$ has length at most
the length of the binary representation of $e$, $b$ and $\mathrm{acc}$ are canonical
binary representations, $b$ has length at most the sum of the lengths of the binary
representations of $a$ and $m$, and $\mathrm{acc}$ has length at most one plus the
length of the binary representation of $m$. If $m>0$, then $b$ and $\mathrm{acc}$ encode
numbers less than $m$, and $\mathrm{acc}\,b^{\mathrm{val}(E)}$ is congruent to $a^e$
modulo $m$.
\end{definition}

\begin{theorem}[Binary length of a remainder]
\lean{Complexity.BinArith.length_bits_mod_le}
\label{Complexity.BinArith.length_bits_mod_le}
\leanok
\uses{Complexity.BinArith.length_bits_le_of_le, Complexity.LogProg.bitsVal}
\difficulty{1}
Let $x$ be a natural number and $w$ a finite bit string whose value is positive. The
binary expansion of the remainder of $x$ upon division by the value of $w$ has length at
most the length of $w$.
\end{theorem}

\begin{theorem}[Initial Power Invariant]
\lean{Complexity.BinArith.powInv_init}
\label{Complexity.BinArith.powInv_init}
\leanok
\uses{Complexity.BinArith.PowInv, Complexity.BinArith.length_bits_bitsVal_le, Complexity.BinArith.length_bits_mod_le, Complexity.BinArith.mod, Complexity.BinArith.powInit, Complexity.LogProg.bitsVal, Complexity.LogProg.bitsVal_bits, Complexity.pairFstD, Complexity.pairFstD_pairEncode, Complexity.pairSndD, Complexity.pairSndD_pairEncode}
\difficulty{5}
For any bit string $z$, let $\pi_1$ and $\pi_2$ be the projections decoded from the
self-delimiting pair encoding, with either projection set to the empty word when
decoding fails, and let $\langle w\rangle$ be the natural number represented by a bit
string with its first bit least significant. The state $\langle B,\langle N,\langle
\operatorname{bits}(A\bmod B),\operatorname{bits}(1\bmod B)\rangle\rangle\rangle$, where
$A=\pi_1(z)$, $N=\pi_1(\pi_2(z))$, and $B=\pi_2(\pi_2(z))$, has the following property:
it can be written as $\langle B,\langle E,\langle b,\mathrm{acc}\rangle\rangle\rangle$
with $|E|\leq |\pi_1(\pi_2(z))|$, $b$ and $\mathrm{acc}$ in canonical binary form,
$|b|\leq |A|+|B|$, and $|\mathrm{acc}|\leq 1+|B|$. If $\langle B\rangle>0$, then
$\langle b\rangle,\langle\mathrm{acc}\rangle<\langle B\rangle$ and
$\langle\mathrm{acc}\rangle\langle b\rangle^{\langle E\rangle}\equiv \langle
A\rangle^{\langle N\rangle}\pmod{\langle B\rangle}$.
\end{theorem}

\begin{theorem}[Preservation and Progress of the Power Invariant]
\lean{Complexity.BinArith.powInv_step}
\label{Complexity.BinArith.powInv_step}
\leanok
\uses{Complexity.BinArith.PowInv, Complexity.BinArith.length_bits_mod_le, Complexity.BinArith.mod, Complexity.BinArith.mul, Complexity.BinArith.powStep, Complexity.LogProg.bitsVal, Complexity.LogProg.bitsVal_bits, Complexity.isNilB, Complexity.pairFstD, Complexity.pairFstD_pairEncode, Complexity.pairSndD, Complexity.pairSndD_pairEncode, Turing.pairEncode}
\difficulty{6}
Let $z$ and $S$ be bit strings, and suppose $S$ has the following property: it is an
encoding $\langle B,\langle E,\langle b,\mathrm{acc}\rangle\rangle\rangle$ with $|E|\leq
|\pi_1(\pi_2(z))|$, with $b$ and $\mathrm{acc}$ in canonical binary form, $|b|\leq
|\pi_1(z)|+|B|$, and $|\mathrm{acc}|\leq 1+|B|$, where $B=\pi_2(\pi_2(z))$; if $\langle
B\rangle>0$, also require $\langle b\rangle,\langle\mathrm{acc}\rangle<\langle B\rangle$
and $\langle\mathrm{acc}\rangle\langle b\rangle^{\langle
E\rangle}\equiv\langle\pi_1(z)\rangle^{\langle\pi_1(\pi_2(z))\rangle}\pmod{\langle
B\rangle}$. Then one step of the specified state update preserves this property.
Moreover, if $\langle B\rangle>0$ and $|\pi_1(\pi_2(S))|>0$, the length of that exponent
component strictly decreases after the update.
\end{theorem}

\begin{theorem}[Length bound from a power witness]
\lean{Complexity.BinArith.length_le_of_powInv}
\label{Complexity.BinArith.length_le_of_powInv}
\leanok
\uses{Complexity.BinArith.PowInv, Complexity.length_pair_components_le, Complexity.pairSndD, Turing.length_pairEncode}
\difficulty{3}
Let $z$ and $S$ be finite bit strings, and suppose there are bit strings $E$, $b$, and
$a$ such that $S$ is the pairing of the total second projection of the total second
projection of $z$ with the nested pair consisting of $E$ and the pair $(b,a)$. Require
that the length of $E$ is at most the length of the total first projection of the total
second projection of $z$, that $b$ and $a$ are the canonical bit strings obtained by
writing their respective values in binary, that the length of $b$ is at most the sum of
the lengths of the total first projection of $z$ and the total second projection of the
total second projection of $z$, and that the length of $a$ is at most one plus the
length of the total second projection of the total second projection of $z$. If the
value of that latter projection is positive, require also that the values of $b$ and $a$
are both smaller than it and that the value of $a$ times the value of $b$ raised to the
value of $E$ is congruent modulo it to the value of the total first projection of $z$
raised to the value of the total first projection of the total second projection of $z$.
Then the length of $S$ is at most $9$ times the sum of the length of $z$ and one.
\end{theorem}

\begin{definition}[Power iteration]
\lean{Complexity.BinArith.powLoop}
\label{Complexity.BinArith.powLoop}
\leanok
\uses{Complexity.BinArith.powInit, Complexity.BinArith.powStep}
For a Boolean list $z$, define its power iteration to be the result of applying the step
function `powStep` exactly $1\cdot(|z|+1)^1$ times, starting from `powInit z`.
\end{definition}

\begin{theorem}[Modular exponentiation is polynomial-time computable]
\lean{Complexity.BinArith.polyTimeComputable_powMod}
\label{Complexity.BinArith.polyTimeComputable_powMod}
\leanok
\uses{Complexity.BinArith.PowInv, Complexity.BinArith.iterate_invariant, Complexity.BinArith.iterate_measure_eq_zero, Complexity.BinArith.length_le_of_powInv, Complexity.BinArith.polyTimeComputable_isZero, Complexity.BinArith.polyTimeComputable_iterate, Complexity.BinArith.polyTimeComputable_mod, Complexity.BinArith.polyTimeComputable_powStep, Complexity.BinArith.polyTimeComputable_slot3, Complexity.BinArith.polyTimeComputable_slot4, Complexity.BinArith.powInit, Complexity.BinArith.powInv_init, Complexity.BinArith.powInv_step, Complexity.BinArith.powLoop, Complexity.BinArith.powMod, Complexity.BinArith.powStep, Complexity.LogProg.bitsVal, Complexity.PolyTimeComputable, Complexity.PolyTimeComputable.comp, Complexity.PolyTimeComputable.pairEncode, Complexity.isNilB, Complexity.length_pairFstD_le, Complexity.length_pairSndD_le, Complexity.pairFstD, Complexity.pairFstD_pairEncode, Complexity.pairSndD, Complexity.pairSndD_pairEncode, Complexity.polyTimeComputable_const, Complexity.polyTimeComputable_ite, Complexity.polyTimeComputable_pairFstD, Complexity.polyTimeComputable_pairSndD}
\difficulty{6}
For a bit string $z$, decode its first component $a$ and the two components $b,m$
obtained by decoding its second component again, using the empty word for any failed
decoding. Interpret each bit string as a binary numeral, with the first bit least
significant, and define the output to be the binary expansion of $0$ if the value of $m$
is $0$, and otherwise of $a^b$ modulo $m$. This function from bit strings to bit strings
is polynomial-time computable.
\end{theorem}

\begin{definition}[Binary greatest common divisor]
\lean{Complexity.BinArith.gcd}
\label{Complexity.BinArith.gcd}
\leanok
\uses{Complexity.LogProg.bitsVal, Complexity.pairFstD, Complexity.pairSndD}
For a Boolean list $z$, define its binary greatest common divisor to be the binary
representation of the greatest common divisor of the natural numbers represented by the
two components decoded from $z$.
\end{definition}

\begin{theorem}[Gcd of encoded binary numbers]
\lean{Complexity.BinArith.gcd_pairEncode}
\label{Complexity.BinArith.gcd_pairEncode}
\leanok
\uses{Complexity.BinArith.gcd, Complexity.LogProg.bitsVal_bits, Turing.pairEncode}
\difficulty{1}
For all natural numbers $a,b$, applying the bit-string gcd operation to the
self-delimiting encoding of the binary expansions of $a$ and $b$ yields the binary
expansion of $\gcd(a,b)$.
\end{theorem}

\begin{definition}[Binary gcd iteration step]
\lean{Complexity.BinArith.gcdStep}
\label{Complexity.BinArith.gcdStep}
\leanok
\uses{Complexity.BinArith.mod, Complexity.LogProg.bitsVal, Complexity.pairFstD, Complexity.pairSndD, Turing.pairEncode}
For a Boolean list $S$, define one step as follows: if the natural number represented by
the second component decoded from $S$ is zero, return $S$; otherwise, return the pair
encoding of that second component and the remainder obtained by applying `mod` to the
pair encoding of the two decoded components.
\end{definition}

\begin{theorem}[Polynomial-Time Computability of One Greatest Common Divisor Step]
\lean{Complexity.BinArith.polyTimeComputable_gcdStep}
\label{Complexity.BinArith.polyTimeComputable_gcdStep}
\leanok
\uses{Complexity.BinArith.gcdStep, Complexity.BinArith.polyTimeComputable_isZero, Complexity.BinArith.polyTimeComputable_mod, Complexity.PolyTimeComputable, Complexity.PolyTimeComputable.comp, Complexity.PolyTimeComputable.pairEncode, Complexity.polyTimeComputable_id, Complexity.polyTimeComputable_ite, Complexity.polyTimeComputable_pairFstD, Complexity.polyTimeComputable_pairSndD}
\difficulty{3}
For a bit string $S$, let $a$ and $b$ be the decoded first and second components, using
the empty string when decoding fails. If $b$ represents zero, return $S$ unchanged;
otherwise return the self-delimiting encoding of $b$ and the binary representation of
$\operatorname{val}(a)\bmod\operatorname{val}(b)$. This function is computable in
polynomial time by a finite Turing machine.
\end{theorem}

\begin{definition}[Gcd invariant]
\lean{Complexity.BinArith.GcdInv}
\label{Complexity.BinArith.GcdInv}
\leanok
\uses{Complexity.LogProg.bitsVal, Complexity.pairFstD, Complexity.pairSndD, Turing.pairEncode}
For Boolean lists $z$ and $S$, the invariant holds when there are Boolean lists $x$ and
$y$ such that $S$ is their pair encoding, each of $x$ and $y$ has length at most the
length of $z$, and the greatest common divisor of the natural numbers represented by $x$
and $y$ equals that of the two components decoded from $z$.
\end{definition}

\begin{definition}[Gcd iteration measure]
\lean{Complexity.BinArith.gcdMeasure}
\label{Complexity.BinArith.gcdMeasure}
\leanok
\uses{Complexity.LogProg.bitsVal, Complexity.pairFstD, Complexity.pairSndD}
For a Boolean list $S$, define its measure as the binary size of the product of the
natural numbers represented by the two components decoded from $S$, plus $1$ if the
first represented number is smaller than the second, and plus $0$ otherwise.
\end{definition}

\begin{theorem}[Strict decrease of binary size]
\lean{Complexity.BinArith.size_lt_of_two_mul_le}
\label{Complexity.BinArith.size_lt_of_two_mul_le}
\leanok
\difficulty{4}
Let $P$ and $P'$ be natural numbers. If $P>0$ and $2P'\le P$, then the binary size of
$P'$ is strictly less than the binary size of $P$.
\end{theorem}

\begin{theorem}[Invariant and decrease under a gcd step]
\lean{Complexity.BinArith.gcdInv_step}
\label{Complexity.BinArith.gcdInv_step}
\leanok
\uses{Complexity.BinArith.GcdInv, Complexity.BinArith.gcdMeasure, Complexity.BinArith.gcdStep, Complexity.BinArith.length_bits_le_of_le, Complexity.BinArith.mod, Complexity.BinArith.size_lt_of_two_mul_le, Complexity.LogProg.bitsVal, Complexity.LogProg.bitsVal_bits, Complexity.pairFstD_pairEncode, Complexity.pairSndD_pairEncode}
\difficulty{6}
Let $z,S$ be finite bit strings, and suppose $S=\langle x,y\rangle$ for bit strings
$x,y$ with $|x|,|y|\le |z|$, such that $\gcd(V(x),V(y))=\gcd(V(\pi_1(z)),V(\pi_2(z)))$.
Here $V$ is the natural number represented by a bit string, and $\pi_1,\pi_2$ are the
total projections of the self-delimiting pair encoding. Define the measure of a state
$S$ as the binary size of $V(\pi_1(S))V(\pi_2(S))$, plus $1$ if
$V(\pi_1(S))<V(\pi_2(S))$ and $0$ otherwise; the gcd step fixes $S$ when its second
projected value is zero and otherwise replaces it by the encoding of that value and the
remainder of the first value on division by the second. The gcd invariant with respect
to $z$ is preserved by this step, the measure strictly decreases whenever it is
positive, and the step fixes $S$ whenever its measure is zero.
\end{theorem}

\begin{definition}[Gcd iteration]
\lean{Complexity.BinArith.gcdLoop}
\label{Complexity.BinArith.gcdLoop}
\leanok
\uses{Complexity.BinArith.gcdStep, Complexity.pairFstD, Complexity.pairSndD, Turing.pairEncode}
For a Boolean list $z$, define its gcd iteration by applying the gcd step exactly
$2(|z|+1)^1$ times, starting from the pair encoding of the two components decoded from
$z$.
\end{definition}

\begin{theorem}[Polynomial-Time Computability of the Greatest Common Divisor]
\lean{Complexity.BinArith.polyTimeComputable_gcd}
\label{Complexity.BinArith.polyTimeComputable_gcd}
\leanok
\uses{Complexity.BinArith.GcdInv, Complexity.BinArith.gcd, Complexity.BinArith.gcdInv_step, Complexity.BinArith.gcdLoop, Complexity.BinArith.gcdMeasure, Complexity.BinArith.gcdStep, Complexity.BinArith.iterate_invariant, Complexity.BinArith.iterate_measure_eq_zero, Complexity.BinArith.polyTimeComputable_gcdStep, Complexity.BinArith.polyTimeComputable_iterate, Complexity.BinArith.polyTimeComputable_normBits, Complexity.LogProg.bitsVal, Complexity.LogProg.bitsVal_lt_two_pow, Complexity.PolyTimeComputable, Complexity.PolyTimeComputable.comp, Complexity.PolyTimeComputable.pairEncode, Complexity.length_pairFstD_le, Complexity.length_pairSndD_le, Complexity.length_pair_components_le, Complexity.pairFstD, Complexity.pairFstD_pairEncode, Complexity.pairSndD, Complexity.pairSndD_pairEncode, Complexity.polyTimeComputable_pairFstD, Complexity.polyTimeComputable_pairSndD, Turing.length_pairEncode, Turing.pairEncode}
\difficulty{6}
For a bit string $z$, decode its two components, taking either component to be empty if
decoding fails, and interpret each as a natural number in least-significant-bit-first
binary. The function returning the binary representation of the greatest common divisor
of those two natural numbers is computable in polynomial time by a finite Turing
machine.
\end{theorem}
